% Chapter 2

\chapter{Problema}	%The main chapter title
%\chaptermark{O \textit{Template}}	%Short version for page header. Comment if not needed
\label{Chapter2}	%For referencing the chapter elsewhere, use \ref{Chapter2} 


\noindent O projeto  de gestão de tarefas pretende demonstrar a tecnologia numa arquitetura distribuída de três camadas. Neste capítulo é apresentada a definição do problema, os requisitos funcionais e técnicos da solução.
%%%%%%%%%%%%%%%%%%%%%%%%%%%%%%%%%%%% Section :Definição

\section{Definição}
\label{sec:Ch2.1}

\noindent No contexto académico, o projeto de gestão de tarefas  é um exemplo criado com o intuito de explicar a tecnologia em estudo. ~O problema surge assim em dois campos distintos:
\begin {itemize}
\item Validação de solução gráfica da tecnologia em estudo.
\item Integração do servidor  \textit{Tomcat} para comunicar com a base de dados  \textit{MySQL} para o serviço remoto da aplicação distribuída.
\end {itemize}

\noindent A solução proposta tem como intuito, validar os conceitos lecionados na \ac{uc} de \ac{rinte}, nomeadamente a comunicação entre um cliente e um servidor através de um serviço \textit{web} baseado no protocolo \ac{http} de três camadas.


\section{Requisitos}
\label{sec:Ch2.2}

\noindent Os requisitos da aplicação Gestão de Tarefas dividem-se em requisitos funcionais, que descrevem o comportamento esperado do sistema do ponto de vista dos seus utilizadores, e requisitos técnicos, que definem as regras não funcionais que a solução deve respeitar.

\subsection{Funcionais}


\noindent A aplicação distribuída tem dois tipos de atores, com níveis de permissão distintos:

\begin{itemize}
	\item \textbf{Administrador} -- responsável pela gestão de todas as tarefas do sistema e pela gestão dos utilizadores registados;
	\item \textbf{Utilizador} -- responsável apenas pela gestão das tarefas que lhe foram atribuídas ou que criou.
\end{itemize}

\noindent A Tabela~\ref{tab:usecases} resume os principais casos de uso identificados, associando-os ao ator responsável e ao método \ac{http} que a interface \ac{rest} do \textit{back-end} segue neste projeto.

\begin{table}[htb]
	\caption{Casos de uso e correspondência com os métodos \ac{http} da \ac{api}.}
	\label{tab:usecases}
	\centering
	\begin{tabular}{l l c}
	\toprule
	\tabhead{Ator} & \tabhead{Caso de Uso} & \tabhead{Método \ac{http}} \\
	\midrule
		Administrador \& Utilizador & Autenticar-se no sistema & \code{POST} \\
		Administrador & Criar tarefa & \code{POST} \\
		Administrador & Consultar todas as tarefas & \code{GET} \\
		Administrador & Atualizar qualquer tarefa & \code{PUT} \\
		Administrador & Eliminar qualquer tarefa & \code{DELETE} \\
		Administrador & Criar utilizador & \code{POST} \\
		Administrador & Consultar utilizadores & \code{GET} \\
		Administrador & Eliminar utilizador & \code{DELETE} \\
		\midrule
		Utilizador & Criar tarefa própria & \code{POST} \\
		Utilizador & Consultar as suas tarefas & \code{GET} \\
		Utilizador & Atualizar tarefa própria & \code{PUT} \\
	\bottomrule
	\end{tabular}
\end{table}

\subsection{Técnicos}

\noindent Para além dos requisitos funcionais, a solução deve ainda respeitar os seguintes requisitos técnicos:

\begin{itemize}
	\item \textbf{Autenticação e autorização} -- o acesso à \ac{api} deve ser condicionado à autenticação do utilizador, devendo o \textit{back-end} validar, em cada pedido, se o ator autenticado possui permissões para a operação solicitada;
	\item \textbf{Portabilidade} -- o cliente gráfico, por ser desenvolvido em \textit{JavaFX}, deve poder ser executado em diferentes sistemas operativos sem alterações ao código;
\end{itemize}
%%%%%%%%%%%%%%%%%%%%%%%%%%%%%%%%%%%%

\section{Sumário}

\noindent Neste capítulo definiu-se o problema, na necessidade de uma aplicação distribuída de gestão de tarefas que integra a tecnologia \textit{JavaFX}, e identificam-se os requisitos funcionais e técnicos que ilustram o desenho do problema.
